\worksheettitle
  {Редактор чисел}
  {\modulemark{M04\(\star\)} Объяснение и конструирование}

\studentnameblock

\begin{warningbox}[1. Две разные ошибки]
  Ученик хотел записать число
  \emph{восемнадцать миллионов сорок тысяч шесть}. Он записал
  \(18\,40\,6\), а прочитал свою запись как «восемнадцать миллионов четыреста
  шесть тысяч».

  \begin{enumerate}[label=\textbf{\arabic*.}]
    \item Запишите задуманное число правильно: \answerblank[48mm].
    \item Объясните ошибку в записи.
      \writinglines{2}
    \item Объясните отдельную ошибку в чтении.
      \writinglines{2}
  \end{enumerate}
\end{warningbox}

\begin{practicebox}[2. Запись, которую можно понять по-разному]
  Запись \(27\,5\,14\) составлена с ошибками. Назовите не менее двух чисел,
  которые автор мог иметь в виду. Для каждого числа запишите возможное чтение.
  \writinglines{4}

  Такая запись называется \textbf{неоднозначной}: без дополнительных сведений
  нельзя определить единственное задуманное число.

  Какое дополнительное условие сделало бы восстановление однозначным?
  \writinglines{1}
\end{practicebox}

\newpage
\section{Конструируем и обобщаем}

\begin{practicebox}[3. Создайте числовую ловушку]
  Придумайте число, в котором:
  \begin{itemize}
    \item один класс полностью нулевой;
    \item другая группа начинается с нуля;
    \item при чтении нулевой класс не произносится.
  \end{itemize}
  Число цифрами: \answerblank[58mm].

  Чтение: \answerblank[105mm].

  Возможная ошибочная запись: \answerblank[58mm].

  Как по ошибке определить, какой класс потерян?
  \writinglines{2}
\end{practicebox}

\begin{practicebox}[4. Составьте пару]
  Запишите два разных числа, которые:
  \begin{itemize}
    \item содержат \(42\) миллиона;
    \item при чтении оканчиваются словом «пять»;
    \item отличаются только группой класса тысяч.
  \end{itemize}
  Числа: \answerblank[45mm] и \answerblank[45mm].

  Прочитайте оба числа.
  \writinglines{2}
\end{practicebox}

\begin{checkpointbox}[5. Обобщите]
  Чем отличаются слово «пять», количество \(5\), группа \(5\) и группа
  \(005\)? Когда допустима каждая запись?
  \writinglines{3}
\end{checkpointbox}
